\begin{claim}{D.3.8}{Proved}{A countable bound for primary infinite-loop presentations of a fixed stable part}
Fix a prime $p$. Let $\{X_i\}_{i\in I}$ be a set or class of
connected infinite-loop spaces with $p$-primary torsion homotopy
groups, at least one of which is noncontractible. Let $n\geq1$
be the least degree in which some $\pi_nX_i$ is nonzero, and put
\[
 \mathcal L=\loc\{\Sigma^\infty X_i:i\in I\},\qquad
 W=\min\mathrm{Rev}\{X_i\to *:i\in I\}.
\]
Then $T(W)=\mathcal L$. If some $\pi_nX_i$ is not divisible,
equivalently some $H_n(X_i;\mathbb F_p)\ne0$, then
\[
 W=\min\mathrm{Rev}
   \{\Omega^\infty E\to *:
                         E\in\mathcal L\cap\Ho(\Sp)_{\geq n}\}.
 \tag{D.3.8a}
\]
If every $\pi_nX_i$ is divisible, then
\[
 W=\min\mathrm{Rev}\left(
   \{\Omega^\infty E\to *:
                   E\in\mathcal L\cap\Ho(\Sp)_{\geq n+1}\}
       \cup\{K(C_{p^\infty},n)\to *\}\right).
 \tag{D.3.8b}
\]
Thus $W$ is determined by $\mathcal L$, the least degree $n$,
and one divisibility flag. For fixed $p$ and $\mathcal L$ there
are at most countably many reversible localizers admitting such
a presentation. If $\mathcal L$ is generated by a small set of
spectra, each is accessible even when the original family $I$
is a class. A family of contractible generators gives the
smallest reversible localizer $W_{\mathbb Z}$ separately.
\end{claim}

\begin{proof}
Choose for each $i$ a connective spectrum $E_i$ giving its indicated
infinite-loop structure, with $X_i\simeq\Omega^\infty E_i$.
Its connective choice has $E_i\in\Ho(\Sp)_{\geq n}$.
The homotopy groups are $p$-primary, so $E_i[1/p]=0$.
The corresponding reduced homology of $X_i$ is $p$-primary as
well, either by the nilpotent localization criterion or by its
Postnikov tower degree by degree.

The exact space--spectrum interface is
\citet[Theorem~2.10 and Proposition~3.1]{Bousfield1996ULP}:
for a space $A$ and a connective spectrum $E$,
\[
 P_A\Omega^\infty E\simeq
          \Omega^\infty P_{\Sigma^\infty A}^{\Sp}E,
 \tag{D.3.8c}
\]
and $A$ kills $\Omega^\infty E$ in ordinary nullification if and
only if $\Sigma^\infty A$ kills $E$ in spectral nullification.
In particular, since $X_i$ kills itself, $E_i$ belongs to the
spectral nullification kernel generated by $\Sigma^\infty X_i$,
and hence to $\mathcal L$. Each $X_i$ is therefore among the
generators on the right of (D.3.8a). The equality
$T(W)=\mathcal L$ also follows at once: $R(\mathcal L)$ contains
every specified aspherical generator, so
$W\subseteq R(\mathcal L)$ and $T(W)\subseteq\mathcal L$;
the suspension spectra of the original generators give the
reverse inclusion.

\medskip\noindent
\emph{The nondivisible case.}
First suppose $I$ is small and let $A=\bigvee_{i\in I}X_i$.
Then $C=\Sigma^\infty A$ is $n$-connective, has $p$-primary
integral homology, and has $H_n(C;\mathbb F_p)\ne0$.
Also $\loc(C)=\mathcal L$.
The spectral conclusion of D.2.9 gives
$P_C^{\Sp}E=0$ for every
$E\in\mathcal L\cap\Ho(\Sp)_{\geq n}$.
Equation (D.3.8c) then makes $\Omega^\infty E$ ordinarily
$A$-acyclic, and consequently $W$-aspherical. This proves
(D.3.8a). The use of ordinary nullification here is justified
by (D.3.8c); it is not an inference from stable asphericity alone.

\medskip\noindent
\emph{The divisible case and the larger GEM-error bound.}
Still suppose $I$ small, and let $A$ and $C$ have the same meaning.
Now $H_n(A;\mathbb Z)=\bigoplus_i\pi_nX_i$ is a nonzero
divisible $p$-primary group. By
\citet[Corollary~3.2]{Bousfield1996ULP}, $A$ kills $SP^\infty A$.
The latter has $K(H_n(A;\mathbb Z),n)$ as a retract, and that
space has $K(C_{p^\infty},n)$ as a retract. Thus this last
space is ordinarily $A$-acyclic. The nonsplit extension
\[
 0\longrightarrow\mathbb F_p\longrightarrow C_{p^\infty}
       \xrightarrow{p}C_{p^\infty}\longrightarrow0
\]
gives a nonzero class in
$H^{n+1}(K(C_{p^\infty},n);\mathbb F_p)$ by the universal
coefficient theorem. Lemma~9.13 of
\citet{EH5Bousfield1997} therefore makes
$K(\mathbb F_p,n+1)$ ordinarily
$K(C_{p^\infty},n)$-acyclic, hence $A$-acyclic.
The precise comparison (D.3.8c), or Proposition~3.1 of its source,
shows that $C$ spectrally kills $\Sigma^{n+1}H\mathbb F_p$.

We claim that
\[
 P_C^{\Sp}E=0
       \quad\text{for every }E\in\mathcal L\cap\Ho(\Sp)_{\geq n+1}.
 \tag{D.3.8d}
\]
Set $C_i=\Sigma^{-i}C$ and $Y_i=P_{C_i}^{\Sp}E$.
The general part of
\citet[Theorem~2.6]{Bousfield1996ULP}, without a support assumption,
says that $F_i=\fib(Y_i\to Y_{i+1})$ is a stable GEM.
Both endpoints are $C_i$-null, so $F_i$ is $C_i$-null.
All $Y_i$ belong to $\mathcal L$, because their nullification
cones are built from shifts of $C$. In particular $F_i$ is
$p$-primary: $\mathcal L$ is contained in the kernel of
inverting $p$.

Since $C_i$ kills $\Sigma^{n+1-i}H\mathbb F_p$, $F_i$ is
null for this latter spectrum. If the stable GEM $F_i$ had a
nonzero homotopy group $G$ in degree $k\geq n+1-i$, its retract
$\Sigma^kHG$ would receive the nonzero group
$\operatorname{Hom}(\mathbb F_p,G)$ in mapping-space degree
$k-(n+1-i)\geq0$. Every nonzero $p$-primary group has an
element of order $p$, so this is impossible. Therefore
\[
 F_i\in\Ho(\Sp)_{\leq n-i},\qquad
 \fib(Y_i\to Y_\infty)\in\Ho(\Sp)_{\leq n-i},
       \quad Y_\infty=\colim_iY_i.
 \tag{D.3.8e}
\]
The second bound follows by extensions and exact filtered
colimits, exactly as in D.2.9.

For a fixed $i$ compare $Y_{i+2}$ with $Y_\infty$.
Their fibre lies in degree at most $n-i-2$, and its suspension
in degree at most $n-i-1$. Since $C_i$ is $(n-i)$-connective,
$t$-structure orthogonality gives
\[
 \Map_{\Sp}(C_i,Y_{i+2})
                 \simeq\Map_{\Sp}(C_i,Y_\infty).
\]
The left side is contractible, since $Y_{i+2}$ is
$C_{i+2}$-null and $C_i=\Sigma^2C_{i+2}$.
Thus $Y_\infty$ is null for every shift of $C$, while it belongs
to $\mathcal L=\loc(C)$, so $Y_\infty=0$.
Equation (D.3.8e) yields $P_C^{\Sp}E\in\Ho(\Sp)_{\leq n}$.
On the other hand nullification by $n$-connective $C$ preserves
$(n+1)$-connectivity of an input which already has it:
the cofiber of any map $\Sigma^kC\to Z$, with $k\geq0$ and
$Z\in\Ho(\Sp)_{\geq n+1}$, still lies in
$\Ho(\Sp)_{\geq n+1}$ by its homotopy exact sequence.
The cone-attachment construction and its transfinite colimits
preserve this bound. Hence $P_C^{\Sp}E=0$, proving (D.3.8d).
The two-step comparison with $Y_{i+2}$, rather than $Y_{i+1}$,
is what controls components with this enlarged error bound.

By (D.3.8c), every higher infinite-loop generator on the right
of (D.3.8b) is ordinarily $A$-acyclic; the Prüfer-space
generator was already shown to be so. This gives one inclusion
in (D.3.8b). For the other write $D_i=\pi_nE_i$.
Each $D_i$ is a direct sum of copies of $C_{p^\infty}$,
possibly zero. Since $K(C_{p^\infty},n)$ is $W$-aspherical,
its suspension spectrum belongs to $\mathcal L$.
Its ordinary acyclic class contains $K(D_i,n)$, and
Proposition~3.1 then puts $\Sigma^nHD_i$ in $\mathcal L$ as well.
Consequently $E_i\in\mathcal L$ implies
$\tau_{\geq n+1}E_i\in\mathcal L$.
In the Postnikov fibre sequence
\[
 \Omega^\infty\tau_{\geq n+1}E_i
          \longrightarrow X_i\longrightarrow K(D_i,n)
\]
both the fibre and the base are aspherical for the right side of
(D.3.8b). The ordinary acyclic class generated by their wedge
is closed under fibre extensions and is contained in that
reversible class. Thus $X_i$ is aspherical there. This uses
fibre and base to obtain the total space; it makes no inference
of fibre acyclicity from the other two terms.

\medskip\noindent
\emph{Class families and the cardinal bound.}
For a class $I$, every individual member of
$\loc\{\Sigma^\infty X_i:i\in I\}$ lies in the localizing
subcategory generated by a small subfamily. Indeed the union
of those subcategories over all small subfamilies is itself
closed under shifts, cofibers and small coproducts: a small
collection of small subfamilies is contained in their small
union. This union contains all the generators and therefore
is $\mathcal L$. Given an individual $E$ in the required
connective range, choose such a small subfamily and adjoin one
source $X_j$ witnessing the least degree and its indicated
nondivisibility or divisibility case. The corresponding
preceding argument applied to its wedge makes
$\Omega^\infty E$ ordinarily acyclic for that wedge, hence
$W$-aspherical. The Prüfer generator and the reverse Postnikov
argument in the divisible case use the same witness. Both
formulas therefore hold in the class case. There are just two
displayed possibilities at each positive integer $n$, proving
the countable bound for fixed $p$ and $\mathcal L$.

Finally suppose $\mathcal L=\loc\{F_a:a\in J\}$ for a small
set $J$. Apply the small-subfamily observation to each $F_a$
and take the union of the resulting source subfamilies,
together with the indicated witness $X_j$. This is a small
subfamily with the same stable part, least degree and
divisibility flag. Its reversible closure equals $W$ by the
appropriate formula, and B.4 makes it accessible.
\end{proof}

This fixes the possible number of primary infinite-loop
presentations of a given stable part. It neither determines
which of the countably many displayed localizers coincide nor
decides the adjacent strictness test in D.2.9. In particular,
it does not replace the nonlinear suspension-nullification
test needed for the zero-finite-part separator.
